\documentclass[t, hyperref={pdfpagelabels=false}]{beamer}
\usetheme{amcg}

\usepackage{listings}
\lstset{basicstyle=\ttfamily}
\usepackage{tikz}
\usepackage{ulem}
\usepackage{hyperref}
\xdefinecolor{lambdablue}{RGB}{19,41,58}
\hypersetup{
  colorlinks=true,
  linkcolor=lambdablue,
  citecolor=lambdablue,
  filecolor=lambdablue,
  urlcolor=lambdablue
}

\title[Parallel Simulations]{Parallel simulations and using clusters}
\subtitle[]{Fluidity Training Course 2015}
\author[]{\footnotesize Simon Mouradian (author: Michael Lange)}
\institute{AMCG, Imperial College London}
\date{\footnotesize 12th November 2015}

\begin{document}
\titlepage

\begin{frame}{Parallel Simulations}
  \begin{itemize}
  \item[] Large simulations require more than a laptop
    \begin{itemize}
      \vspace{1em}
    \item Fluidity runs in parallel via MPI
    \item Optional: OpenMP threading for hybrid runs
      \begin{itemize}
      \item Configuration option \lstinline+--enable-openmp+
      \end{itemize}

      \vspace{1em}
    \item Fluidity runs on various clusters
      \begin{itemize}
      \item Imperial College: CX1, CX2 (Intel)
      \item UK national supercomputer: Archer (Cray CX30)
      \item ANU/NCI: Raijin (Fujitsu Primergy)
      \end{itemize}
    \end{itemize}
  \end{itemize}
\end{frame}

\begin{frame}{Parallel Simulations}
  \begin{columns}[t]
    \column[T]{0.7\textwidth}

    \begin{itemize}
    \item[] How to run Fluidity in parallel:
      \begin{itemize}
        \vspace{1em}
      \item No changes to the .flml required!
        \begin{itemize}
        \item Parallel compatible preconditioner (not eisenstat)
        \item Optional: Remove fields from output (.stat and .vtu)
        \end{itemize}

        \vspace{1em}
      \item Need to decompose the mesh
        \begin{itemize}
        \item Tools provided:
          \\\lstinline+flredecomp+, \sout{\lstinline+fldecomp+}
          \vspace{2pt}
        \item Build both with: \lstinline+make fltools+
        \end{itemize}
      \end{itemize}
    \end{itemize}

    \column[T]{0.3\textwidth}
    \vspace{2em}
    \begin{figure}[ht] \centering
      \begin{tikzpicture}[scale=0.8]
        \tikzstyle{line} = [draw, -, thick]
\tikzstyle{nodraw} = [draw, fill, circle, minimum width=0pt, inner sep=0pt]
\tikzstyle{nocolour} = [draw=black!0, fill=black!0]
\tikzstyle{dark} = [draw=black, fill=black]
\tikzstyle{darkred} = [draw=black!20!red, fill=black!20!red]
\tikzstyle{darkblue} = [draw=black!45!blue, fill=black!45!blue]

\tikzstyle{vsmall} = [circle, minimum width=5pt, inner sep=0pt]
\tikzstyle{vint} = [dark]
\tikzstyle{vext} = [dark]
\tikzstyle{fext} = [dark]
\tikzstyle{fint} = [dark]
\tikzstyle{rcorefill} = [nocolour]
\tikzstyle{rownedfill} = [nocolour]
\tikzstyle{rhalofill} = [nocolour]
\tikzstyle{lpart} = [darkred]
\tikzstyle{rcore} = [darkblue]
\tikzstyle{rowned} = [darkblue]
\tikzstyle{rhalo} = [darkred]

\begin{scope}[overlay, remember picture]
\filldraw[rcorefill] (1.9,0.8) -- (2.2, 1.7) -- (2.3, 2.5) -- (1.6,3.3) -- 
                     (1.5,2.3) -- (1.4,1.6) -- cycle;
\end{scope}

\begin{scope}[overlay, remember picture]
\filldraw[rownedfill] (1.9,0.8) -- (1.4,1.6) -- (1.5,2.3) -- (1.6,3.3) --
                      (0.3,3.6) -- (0.7,2.9) -- (0.6,2.) -- (0.9,1.1) --
                      (1.1,0.4) -- cycle;
\end{scope}

\begin{scope}[overlay, remember picture]
\filldraw[rhalofill] (0.3,3.6) -- (0.7,2.9) -- (0.6,2.) -- (0.9,1.1) --
                     (1.1,0.4) -- (.1,0) -- (.1,0.9) -- (-0.1,1.8) -- (-0.2,2.7) --
                     (-1.,3.2) -- cycle;


\end{scope}

%% Boundary vertices
\node (0) [vsmall, vext, lpart, rhalo] at (.1,0) {};
\node (1) [vsmall, vext, rowned] at (1.1,0.4) {};
\node (2) [vsmall, vext, rowned] at (1.9,0.8) {};
\node (3) [vsmall, vext, rcore] at (2.2, 1.7) {} ;
\node (4) [vsmall, vext, rcore] at (2.3, 2.5) {};
\node (5) [vsmall, vext, rowned] at (1.6,3.3) {};
\node (6) [vsmall, vext, rowned] at (0.3,3.6) {};
\node (7) [vsmall, vext, lpart, rhalo] at (-1.,3.2) {};
\node (8) [vsmall, vext, lpart] at (-1.8,2.4) {};
\node (9) [vsmall, vext, lpart] at (-1.6,1.3) {};
\node (10) [vsmall, vext, lpart] at (-0.9,0.4) {};

%% Boundary facets
\draw[line, fext, rhalo] (0) -- (1);
\draw[line, fext, rowned] (1) -- (2);
\draw[line, fext, rcore] (2) -- (3);
\draw[line, fext, rcore] (3) -- (4);
\draw[line, fext, rcore] (4) -- (5);
\draw[line, fext, rowned] (5) -- (6);
\draw[line, fext, rhalo] (6) -- (7);
\draw[line, fext, lpart] (7) -- (8);
\draw[line, fext, lpart] (8) -- (9);
\draw[line, fext, lpart] (9) -- (10);
\draw[line, fext, lpart] (10) -- (0);


%% Interior vertices, right partition
\node (20) [vsmall, vint, rowned] at (0.9,1.1) {};
\node (21) [vsmall, vint, rowned] at (1.4,1.6) {};
\node (22) [vsmall, vint, rowned] at (1.5,2.3) {};
\node (23) [vsmall, vint, rowned] at (0.7,2.9) {};
\node (24) [vsmall, vint, rowned] at (0.6,2.) {};

%% Interior vertices, left partition
\node (30) [vsmall, vint, lpart, rhalo] at (.1,0.9) {};
\node (31) [vsmall, vint, lpart, rhalo] at (-0.1,1.8) {};
\node (32) [vsmall, vint, lpart, rhalo] at (-0.2,2.7) {};
\node (33) [vsmall, vint, lpart] at (-1.,2.2) {};
\node (34) [vsmall, vint, lpart] at (-0.7,1.1) {};


%% Interior facets
\draw[line, fint, rowned] (20) -- (1);
\draw[line, fint, rowned] (20) -- (2);
\draw[line, fint, rowned] (20) -- (21);
\draw[line, fint, rowned] (20) -- (24);
\draw[line, fint, rhalo] (20) -- (30);
\draw[line, fint, rhalo] (20) -- (31);

\draw[line, fint, rowned] (21) -- (2);
\draw[line, fint, rcore] (21) -- (3);
\draw[line, fint, rowned] (21) -- (22);
\draw[line, fint, rowned] (21) -- (24);

\draw[line, fint, rcore] (22) -- (3);
\draw[line, fint, rcore] (22) -- (4);
\draw[line, fint, rowned] (22) -- (5);
\draw[line, fint, rowned] (22) -- (23);
\draw[line, fint, rowned] (22) -- (24);

\draw[line, fint, rowned] (23) -- (5);
\draw[line, fint, rowned] (23) -- (6);
\draw[line, fint, rowned] (23) -- (24);
\draw[line, fint, rhalo] (23) -- (32);

\draw[line, fint, rhalo] (24) -- (31);
\draw[line, fint, rhalo] (24) -- (32);

\draw[line, fint, lpart] (30) -- (10);
\draw[line, fint, lpart, rhalo] (30) -- (0);
\draw[line, fint, rhalo] (30) -- (1);
\draw[line, fint, lpart, rhalo] (30) -- (31);
\draw[line, fint, lpart] (30) -- (34);

\draw[line, fint, lpart, rhalo] (31) -- (32);
\draw[line, fint, lpart] (31) -- (33);
\draw[line, fint, lpart] (31) -- (34);

\draw[line, fint, rhalo] (32) -- (6);
\draw[line, fint, lpart, rhalo] (32) -- (7);
\draw[line, fint, lpart] (32) -- (33);

\draw[line, fint, lpart] (33) -- (7);
\draw[line, fint, lpart] (33) -- (8);
\draw[line, fint, lpart] (33) -- (9);
\draw[line, fint, lpart] (33) -- (34);

\draw[line, fint, lpart] (34) -- (9);
\draw[line, fint, lpart] (34) -- (10);

      \end{tikzpicture}
    \end{figure}
  \end{columns}
\end{frame}

\begin{frame}{Mesh decomposition}
  \begin{itemize}
  \item[] Mesh decomposition before the run:
    \begin{itemize}

      \vspace{1em}
    \item \lstinline+mpiexec -n 4 flredecomp -i 1 -o 4 serial parallel+
      \begin{itemize}
      \item Needs to be run in parallel!
      \item Decomposes mesh named in \lstinline+serial.flml+
      \item Creates new parallel setup: \lstinline+parallel.flml+
      \end{itemize}

      \vspace{1em}
    \item Now run this on 4 processors: \\
      \lstinline+mpiexec -n 4 fluidity parallel.flml+
    \end{itemize}
  \end{itemize}
\end{frame}

\begin{frame}{Mesh decomposition}
  \begin{itemize}
  \item[] Mesh decomposition with checkpoints:
    \begin{itemize}
      \vspace{1em}
    \item Switch on checkpointing!
      \begin{itemize}
      \item \lstinline+/io/checkpointing+ in .flml
      \end{itemize}

      \vspace{1em}
    \item \lstinline+mpiexec -n 8 flredecomp -i 4 -o 8 par_4 par_8+
      \begin{itemize}
      \item Needs to be run on \lstinline+max(-i, -o)+ processes
      \item Enables re-starting on more/fewer processes
      \end{itemize}

      \vspace{1em}
    \item Run as before: \\
      \lstinline+mpiexec -n 8 fluidity par_8.flml+
    \end{itemize}
  \end{itemize}
\end{frame}

\begin{frame}{Running Fluidity on clusters}
  \begin{itemize}
  \item Set up simulation on local machine!
    \vspace{1em}
  \item Install Fluidity on cluster
    \begin{itemize}
    \item CX1/CX2: \lstinline+module load fluidity+
    \item Archer: \href{https://github.com/FluidityProject/buildscripts/tree/master/uk/ac/archer}{instructions here}
    \end{itemize}
  \item Create \lstinline+.pbs+ submission script
    \begin{itemize}
    \item Submit with \lstinline+qsub+
    \end{itemize}
  \item Copy results back to local machine
    \begin{itemize}
    \item Visualisation as before
    \end{itemize}
  \end{itemize}
\end{frame}

\begin{frame}[fragile]{CX1}
  \begin{lstlisting}[language=bash, basicstyle=\tiny]
#!/bin/bash
#PBS -N pile
#PBS -l walltime=0:20:00
# Resource requirements
#PBS -l select=1:ncpus=8:mem=1550mb

module load fluidity

echo Working directory is $PBS_O_WORKDIR
cd $PBS_O_WORKDIR

mpiexec flredecomp -i 4 -o 8 flow_past_pile_50_checkpoint flow_past_pile_50_cp_n8

mpiexec fluidity -v2 -l flow_past_pile_50_cp_n8.flml
  \end{lstlisting}
\end{frame}

\begin{frame}[fragile]{Archer}
  \begin{lstlisting}[language=bash, basicstyle=\tiny]
#!/bin/bash --login
#PBS -l walltime=01:00:00
#PBS -N jet_sim
#PBS -l select=2
#PBS -A n01-IC1

# Load Fluidity environment
module swap PrgEnv-cray PrgEnv-fluidity

# Make sure any symbolic links are resolved to absolute path
export PBS_O_WORKDIR=$(readlink -f $PBS_O_WORKDIR)
cd $PBS_O_WORKDIR

# Workaround for python environment
export LD_LIBRARY_PATH=$ANACONDA_LIB:$LD_LIBRARY_PATH

export FLUIDITYDIR=path/to/my/fluidity/install

# Decompose the mesh
aprun -n 48 -N 24 -d 1 -j 1 $FLUIDITYDIR/bin/flredecomp -i 1 -o 48 jet_sim set_sim_48

# Run simulation
aprun -n 48 -N 24 -d 1 -j 1 $FLUIDITYDIR/bin/fluidity -l -p -v2 jet_sim_48.flml
  \end{lstlisting}
\end{frame}

\end{document}
